% GENERATED FILE. Edit canonical lessons, then run npm run generate:books.
% canonical-source-hash: fnv1a64:42e8ceab90e45e75
% canonical-lessons: SA-C187-dayaluh, SA-C187-udarah, SA-C187-krpanah, SA-C187-satyavadi, SA-C187-patuh

\chapter{Kind, Generous, Stingy, Truthful, Adept}
\label{ch:sa-kind-generous-stingy-truthful-adept}
\hlchaptermodality{187}

\begin{chapteropening}
\textbf{By the end of this chapter:} \emph{I can say kind, generous, stingy, truthful and adept.}

Complete the last lesson of chapter 187: I can say kind, generous, stingy, truthful and adept.
\end{chapteropening}

\section[\texorpdfstring{\sk{दयालुः} (dayāluḥ)}{दयालुः (dayāluḥ)}]{\sk{दयालुः} (dayāluḥ) — kind}

\label{lesson:SA-C187-dayaluh}



\begin{quote}
Before the new one: say the Sanskrit for fine, then the Sanskrit for steady.
\end{quote}

\subsection*{You'll want to know: \sk{दयालुः}}
\textbf{\sk{दयालुः}} — \emph{dayāluḥ} — \textquotedblleft{}kind\textquotedblright{}.

\begin{grammarlens}[title={What you've built}]
One more word for this chapter.
\end{grammarlens}

\subsection*{Your turn}
\begin{itemize}
\raggedright
  \item \emph{Say it:} \emph{dayāluḥ}
  \item \emph{Say it:} \emph{dayāluḥ}, once more
  \item \emph{Say it:} say \emph{dayāluḥ}
  \item \emph{Recall:} say the Sanskrit for fine, then the Sanskrit for steady, then say \emph{dayāluḥ} again
\end{itemize}

\begin{tcolorbox}[breakable,colback=teal!4,colframe=teal!35!black,title={Before you move on}]
What does \sk{दयालुः} mean? (\textquotedblleft{}Kind\textquotedblright{}.) Say it once more.
\end{tcolorbox}
\section[\texorpdfstring{\sk{उदारः} (udāraḥ)}{उदारः (udāraḥ)}]{\sk{उदारः} (udāraḥ) — generous}

\label{lesson:SA-C187-udarah}



\begin{quote}
Before the new one: say the Sanskrit for steady, then the Sanskrit for kind.
\end{quote}

\subsection*{You'll want to know: \sk{उदारः}}
\textbf{\sk{उदारः}} — \emph{udāraḥ} — \textquotedblleft{}generous\textquotedblright{}.

\begin{grammarlens}[title={What you've built}]
One more word for this chapter.
\end{grammarlens}

\subsection*{Your turn}
\begin{itemize}
\raggedright
  \item \emph{Say it:} \emph{udāraḥ}
  \item \emph{Say it:} \emph{udāraḥ}, once more
  \item \emph{Say it:} say \emph{udāraḥ}
  \item \emph{Recall:} say the Sanskrit for steady, then the Sanskrit for kind, then say \emph{udāraḥ} again
\end{itemize}

\begin{tcolorbox}[breakable,colback=teal!4,colframe=teal!35!black,title={Before you move on}]
What does \sk{उदारः} mean? (\textquotedblleft{}Generous\textquotedblright{}.) Say it once more.
\end{tcolorbox}
\section[\texorpdfstring{\sk{कृपणः} (kṛpaṇaḥ)}{कृपणः (kṛpaṇaḥ)}]{\sk{कृपणः} (kṛpaṇaḥ) — stingy}

\label{lesson:SA-C187-krpanah}



\begin{quote}
Before the new one: say the Sanskrit for kind, then the Sanskrit for generous.
\end{quote}

\subsection*{You'll want to know: \sk{कृपणः}}
\textbf{\sk{कृपणः}} — \emph{kṛpaṇaḥ} — \textquotedblleft{}stingy\textquotedblright{}.

\begin{grammarlens}[title={What you've built}]
One more word for this chapter.
\end{grammarlens}

\subsection*{Your turn}
\begin{itemize}
\raggedright
  \item \emph{Say it:} \emph{kṛpaṇaḥ}
  \item \emph{Say it:} \emph{kṛpaṇaḥ}, once more
  \item \emph{Say it:} say \emph{kṛpaṇaḥ}
  \item \emph{Recall:} say the Sanskrit for kind, then the Sanskrit for generous, then say \emph{kṛpaṇaḥ} again
\end{itemize}

\begin{tcolorbox}[breakable,colback=teal!4,colframe=teal!35!black,title={Before you move on}]
What does \sk{कृपणः} mean? (\textquotedblleft{}Stingy\textquotedblright{}.) Say it once more.
\end{tcolorbox}
\section[\texorpdfstring{\sk{सत्यवादी} (satyavādī)}{सत्यवादी (satyavādī)}]{\sk{सत्यवादी} (satyavādī) — truthful}

\label{lesson:SA-C187-satyavadi}



\begin{quote}
Before the new one: say the Sanskrit for generous, then the Sanskrit for stingy.
\end{quote}

\subsection*{You'll want to know: \sk{सत्यवादी}}
\textbf{\sk{सत्यवादी}} — \emph{satyavādī} — \textquotedblleft{}truthful\textquotedblright{}.

\begin{grammarlens}[title={What you've built}]
One more word for this chapter.
\end{grammarlens}

\subsection*{Your turn}
\begin{itemize}
\raggedright
  \item \emph{Say it:} \emph{satyavādī}
  \item \emph{Say it:} \emph{satyavādī}, once more
  \item \emph{Say it:} say \emph{satyavādī}
  \item \emph{Recall:} say the Sanskrit for generous, then the Sanskrit for stingy, then say \emph{satyavādī} again
\end{itemize}

\begin{tcolorbox}[breakable,colback=teal!4,colframe=teal!35!black,title={Before you move on}]
What does \sk{सत्यवादी} mean? (\textquotedblleft{}Truthful\textquotedblright{}.) Say it once more.
\end{tcolorbox}
\section[\texorpdfstring{\sk{पटुः} (paṭuḥ)}{पटुः (paṭuḥ)}]{\sk{पटुः} (paṭuḥ) — adept, skilful}

\label{lesson:SA-C187-patuh}



\begin{quote}
Before the new one: say the Sanskrit for stingy, then the Sanskrit for truthful.
\end{quote}

\subsection*{You'll want to know: \sk{पटुः}}
\textbf{\sk{पटुः}} — \emph{paṭuḥ} — \textquotedblleft{}adept, skilful\textquotedblright{}.

\begin{grammarlens}[title={What you've built}]
One more word for this chapter.
\end{grammarlens}

\subsection*{Your turn}
\begin{itemize}
\raggedright
  \item \emph{Say it:} \emph{paṭuḥ}
  \item \emph{Say it:} \emph{paṭuḥ}, once more
  \item \emph{Say it:} say \emph{paṭuḥ}
  \item \emph{Recall:} say the Sanskrit for stingy, then the Sanskrit for truthful, then say \emph{paṭuḥ} again
\end{itemize}

\begin{tcolorbox}[breakable,colback=teal!4,colframe=teal!35!black,title={Before you move on}]
What does \sk{पटुः} mean? (\textquotedblleft{}Adept, skilful\textquotedblright{}.) Say it once more.
\end{tcolorbox}
